%=====================================================================
\chapter{Data Science for Cybersecurity}\label{ch:cybersecurity}
%=====================================================================
\locator{5}{Part V --- Data Science in Systems Context}

\begin{outcomes}
By the end of this chapter you will be able to:
\begin{tight}
  \item \textbf{Identify} the principal security data sources and what each can and cannot reveal.
  \item \textbf{Design} detection for intrusion, malware, phishing and insider misuse, choosing supervised or unsupervised methods appropriately.
  \item \textbf{Quantify} the alert-fatigue problem and set thresholds that a real SOC can sustain.
  \item \textbf{Explain} the four classes of attack against machine learning systems themselves.
  \item \textbf{Specify} defences for a deployed security model against an adaptive adversary.
\end{tight}
\end{outcomes}

\begin{prereq}
\chref{probability} (base rates --- the arithmetic that governs every alerting
decision), \chref{evaluation} (imbalance, PR curves, cost-based thresholds),
\chref{unsupervised} (anomaly detection) and \chref{timeseries} (streaming and
drift).
\end{prereq}

\begin{keyterms}
SIEM $\bullet$ NetFlow $\bullet$ intrusion detection $\bullet$ signature vs anomaly
detection $\bullet$ UEBA $\bullet$ baseline $\bullet$ alert fatigue $\bullet$
threat intelligence $\bullet$ adversarial machine learning $\bullet$ evasion
$\bullet$ poisoning $\bullet$ model extraction $\bullet$ membership inference
$\bullet$ defence in depth
\end{keyterms}

\section{What Makes This Domain Different}

\begin{alertbox}[The Assumption That Fails Here]
Every method in Part III assumes the data is generated by a process that does not
care what your model does. In security that assumption is false. There is a human
adversary who will observe your defence, work out its boundary and step around it.
This changes the discipline: a security model is not fitted once and monitored ---
it is \emph{contested continuously}. Nothing else in this book faces that.
\end{alertbox}

\section{Security Data Sources}

\rowcolors{2}{white}{lightbg}
\begin{longtable}{@{}p{2.9cm}p{4.6cm}p{6.7cm}@{}}
\toprule
\rowcolor{primary!15}
\textbf{Source} & \textbf{Contains} & \textbf{Blind spot you must know about} \\
\midrule
\endhead
Authentication logs & Who logged in, from where, success or failure & Says nothing about what they did afterwards \\
NetFlow / traffic metadata & Who talked to whom, how much, for how long & No payload; encrypted traffic looks identical to benign \\
Endpoint (EDR) & Processes, files, registry, command lines & Only on hosts with an agent installed --- rarely all of them \\
DNS logs & Every name resolved & Very high volume; heavy encoding of exfiltration \\
Firewall / proxy & Allowed and blocked connections & Only what traverses the boundary; lateral movement is invisible \\
Application logs & Business-level actions & Format varies per app; often the first thing an attacker disables \\
Threat intelligence & Known-bad indicators & Only covers what someone has already seen and published \\
\bottomrule
\end{longtable}

\begin{conceptbox}[Signature versus Anomaly Detection]
\term{Signature} detection matches known-bad patterns: precise, fast, explainable,
and structurally incapable of catching anything new. \term{Anomaly} detection flags
deviation from normal: can catch novel attacks, and generates far more false
positives. Real systems run both --- signatures handle the known volume cheaply so
that analysts' attention is available for the anomalies.
\end{conceptbox}

\section{The Alert Fatigue Problem}

\begin{worked}[Sizing a SOC's Alert Budget]
An organisation processes 2 million authentication events daily. Roughly 20 are
genuinely malicious ($10^{-5}$ base rate). Three analysts each investigate about
40 alerts a day, so the \textbf{sustainable budget is 120 alerts/day}.

\smallskip
\textbf{Detector A --- 95\% recall, 0.1\% false-positive rate.}
\begin{tight}
  \item True positives: $20 \times 0.95 = 19$
  \item False positives: $1{,}999{,}980 \times 0.001 = 2{,}000$
  \item Alerts/day: $2{,}019$ --- \textbf{17$\times$ over budget}
  \item Precision: $19/2019 = 0.9\%$
\end{tight}
Unworkable. Within two weeks analysts triage by gut feel and the tool is
effectively off.

\smallskip
\textbf{Detector B --- 70\% recall, 0.005\% false-positive rate.}
\begin{tight}
  \item True positives: $14$
  \item False positives: $1{,}999{,}986 \times 0.00005 = 100$
  \item Alerts/day: $114$ --- \textbf{within budget}
  \item Precision: $14/114 = 12.3\%$
\end{tight}

\smallskip
\textbf{B is the better detector} despite catching five fewer real attacks in
theory, because A's alerts will not actually be investigated. The realistic
comparison is 14 attacks investigated versus perhaps 3 --- whichever of A's 2{,}019
alerts happen to get looked at.

\smallskip
\textbf{The design rule:} \emph{start from the analyst capacity and work backwards
to the required false-positive rate.} Here, 120 alerts/day against 2 million events
demands an FPR below $6\times10^{-5}$. Any model that cannot reach that is not
deployable at this volume, however good its ROC-AUC --- which, note, would be
excellent for both detectors.
\end{worked}

\dsfig{%
\begin{tikzpicture}
\begin{axis}[
  width=12.4cm, height=5.4cm,
  xmode=log, xmin=1e-6, xmax=1e-2,
  ymin=0, ymax=2600,
  xlabel={\scriptsize false-positive rate (log scale)},
  ylabel={\scriptsize alerts per day},
  axis lines=left, grid=major, grid style={gray!18},
  tick label style={font=\tiny}, label style={font=\tiny},
  legend style={font=\tiny, draw=gray!40, at={(0.02,0.98)}, anchor=north west},
]
\addplot[accent, line width=1.4pt, domain=1e-6:1e-2, samples=120] {2000000*x + 15};
\addlegendentry{alerts generated}
\addplot[greenhl, dashed, line width=1.2pt, domain=1e-6:1e-2] {120};
\addlegendentry{analyst capacity (3 staff)}

\addplot[accent, only marks, mark=*, mark size=2.2pt] coordinates {(0.001,2019)};
\node[font=\tiny, accent, anchor=east] at (axis cs:0.0009,2100) {Detector A};
\addplot[primary, only marks, mark=*, mark size=2.2pt] coordinates {(0.00005,114)};
\node[font=\tiny, primary, anchor=west, align=left] at (axis cs:0.00006,330)
     {Detector B\\(deployable)};

\draw[gray!55, dashed] (axis cs:0.00006,0) -- (axis cs:0.00006,120);
\node[font=\tiny, text=gray!55!black, anchor=west, align=left, rotate=90]
     at (axis cs:0.000065,300) {max viable FPR};
\end{axis}
\end{tikzpicture}}%
{Alert volume against false-positive rate at 2 million events per day. The relationship is
linear and brutal: the analyst capacity line fixes a maximum FPR, and that constraint ---
not accuracy --- determines whether a detector can be deployed.}{alert-fatigue}

\section{Four Detection Problems}

\dsfig{%
\begin{tikzpicture}[font=\scriptsize,
  hd/.style={rounded corners=3pt, fill=#1, text=white, font=\scriptsize\bfseries,
             text width=6.7cm, align=center, minimum height=0.55cm},
  b/.style={rounded corners=3pt, draw=#1, fill=#1!7, text width=6.7cm, align=left,
            font=\scriptsize, inner sep=6pt, minimum height=2.5cm}]

\node[hd=secondary] at (0,0)    {INTRUSION DETECTION};
\node[b=secondary]  at (0,-1.62){%
  \textbf{Data:} NetFlow, firewall, endpoint\\
  \textbf{Features:} bytes in/out ratio, distinct ports per host per minute, session
  duration, deviation from the host's own 30-day profile\\
  \textbf{Method:} gradient boosting where labels exist; isolation forest where they do
  not\\
  \textbf{Hard part:} encrypted traffic hides payload, so only metadata is available};

\node[hd=goldhl] at (7.3,0)    {MALWARE CLASSIFICATION};
\node[b=goldhl]  at (7.3,-1.62){%
  \textbf{Data:} binaries, sandbox traces\\
  \textbf{Features:} imported API calls, byte histograms, string entropy, section sizes;
  or raw bytes as an image (\chref{architectures})\\
  \textbf{Method:} tree ensembles or CNNs\\
  \textbf{Hard part:} packing and polymorphism mean the same malware looks different every
  time};

\node[hd=purplehl] at (0,-3.55)    {PHISHING DETECTION};
\node[b=purplehl]  at (0,-5.17){%
  \textbf{Data:} email text, headers, URLs\\
  \textbf{Features:} TF-IDF (\chref{nlp}), sender reputation, URL/domain age, homoglyphs,
  reply-to mismatch\\
  \textbf{Method:} naive Bayes or logistic regression --- fast enough for gateway volumes\\
  \textbf{Hard part:} LLM-written phishing removed the spelling-error signal};

\node[hd=tealhl] at (7.3,-3.55)    {UEBA --- INSIDER MISUSE};
\node[b=tealhl]  at (7.3,-5.17){%
  \textbf{Data:} auth, file access, application logs\\
  \textbf{Features:} deviation from \emph{this user's own} baseline --- hours, volume,
  systems touched\\
  \textbf{Method:} per-entity anomaly scoring\\
  \textbf{Hard part:} legitimate role changes look identical to compromise; monitoring
  staff raises real ethical questions};
\end{tikzpicture}}%
{Four detection problems, their data, their standard methods and the difficulty that
defines each. Note that three of the four rely on \emph{per-entity baselines} rather than
global thresholds.}{security-problems}

\begin{conceptbox}[The Single Most Useful Feature Idea in Security Analytics]
Compare an entity against \emph{its own history}, not against the population. A
server that suddenly contacts 200 external addresses is alarming even if 200 is
unremarkable for other servers; a user working at 3 a.m.\ matters only if that user
never has. Per-entity baselines convert an intractable global-threshold problem
into a tractable per-entity one, and they are what makes UEBA work.
\end{conceptbox}

\section{Attacking the Model Itself}

\begin{definitionbox}[Adversarial Machine Learning]
Attacks that target the ML system rather than the network around it. Four classes,
distinguished by \emph{when} they act and \emph{what} they want.
\end{definitionbox}

\dsfig{%
\begin{tikzpicture}[font=\scriptsize,
  hd/.style={rounded corners=3pt, fill=accent, text=white, font=\scriptsize\bfseries,
             text width=6.7cm, align=center, minimum height=0.55cm},
  b/.style={rounded corners=3pt, draw=accent, fill=accent!6, text width=6.7cm, align=left,
            font=\scriptsize, inner sep=6pt, minimum height=2.35cm}]

\node[hd] at (0,0)    {1. EVASION (at inference)};
\node[b]  at (0,-1.55){%
  \textbf{Goal:} be classified as benign.\\[3pt]
  Perturb the input just enough to cross the boundary --- pad a binary, insert
  invisible characters, throttle a scan below the rate threshold.\\[3pt]
  {\color{greenhl!45!black}\textbf{Defences:} adversarial training; ensembles;
   features costly for the attacker to change; randomised thresholds}};

\node[hd] at (7.3,0)    {2. POISONING (at training)};
\node[b]  at (7.3,-1.55){%
  \textbf{Goal:} corrupt the next model.\\[3pt]
  Feed crafted traffic during the data collection window so the retrained model
  learns to treat the attack pattern as normal.\\[3pt]
  {\color{greenhl!45!black}\textbf{Defences:} vet training data; anomaly-screen the
   training set; human review of label changes; do not auto-retrain on unvetted
   production data}};

\node[hd] at (0,-3.45)    {3. MODEL EXTRACTION};
\node[b]  at (0,-5.0){%
  \textbf{Goal:} copy your model.\\[3pt]
  Query repeatedly and train a substitute on the answers, then attack the substitute
  offline until evasion is reliable.\\[3pt]
  {\color{greenhl!45!black}\textbf{Defences:} rate limiting; return decisions rather
   than scores; monitor for systematic probing patterns}};

\node[hd] at (7.3,-3.45)    {4. MEMBERSHIP INFERENCE};
\node[b]  at (7.3,-5.0){%
  \textbf{Goal:} learn who was in the training data.\\[3pt]
  Exploit the model's higher confidence on examples it was trained on --- a privacy
  breach, not an evasion.\\[3pt]
  {\color{greenhl!45!black}\textbf{Defences:} regularisation; differential privacy;
   do not expose raw confidence scores}};

\node[anchor=north west, align=left, text width=13.4cm, font=\scriptsize,
      text=primary!60!black] at (-3.35,-6.75)
     {\textbf{Note the pattern in the defences.} Almost none of them is a better model.
      They are rate limits, data vetting, output restriction and human review --- process
      and architecture controls. Security of an ML system is not primarily achieved by
      improving its accuracy, exactly as agent safety in \chref{agentic} was not achieved
      by improving its prompt.};
\end{tikzpicture}}%
{The four classes of attack on machine learning systems. Evasion and extraction happen at
inference time; poisoning happens at training time; membership inference targets privacy
rather than performance.}{adversarial-ml}

\begin{alertbox}[Retraining on Production Data Is an Attack Surface]
The advice from \chref{mlops} --- retrain regularly on recent data to counter drift
--- is exactly the mechanism a poisoning attack exploits. In security you cannot
have both automatic retraining and unvetted training data. Either screen the
training set for anomalies before use, or keep a human in the approval path for
model promotion. This is a genuine tension between two pieces of good advice, and
the resolution is domain-specific.
\end{alertbox}

\begin{worked}[A Concrete Evasion, and Why Feature Choice Is a Defence]
A detector flags port scans using \texttt{distinct\_ports\_per\_minute > 50}.

\smallskip
\textbf{The evasion:} the attacker scans 40 ports per minute instead. Detection
drops to zero; the scan takes 25\% longer, which costs the attacker almost nothing.

\textbf{Why the model was fragile:} the feature is \emph{cheap for the attacker to
control}. Any threshold on a directly-controllable quantity is one experiment away
from being defeated.

\smallskip
\textbf{A more robust feature set:}
\begin{tight}
  \item \texttt{distinct\_ports\_per\_hour} --- a longer window; slowing further
        costs real time.
  \item \texttt{fraction\_of\_ports\_that\_are\_closed} --- a scanner necessarily
        hits closed ports; legitimate traffic rarely does. Hard to avoid without
        already knowing the answer.
  \item \texttt{deviation from this source's own 30-day profile} --- a host that has
        never initiated outbound connections doing so at all is anomalous at any
        rate.
  \item \texttt{sequential-port-ordering score} --- avoidable, but only by
        randomising, which is itself detectable.
\end{tight}

\smallskip
\textbf{The principle:} prefer features that are \emph{expensive or self-defeating
for the attacker to manipulate}. In ordinary machine learning you choose features
for predictive power; in security you choose them for predictive power \emph{and}
manipulation cost. That is a genuinely different design criterion, and it is the
main thing this chapter adds to \chref{features}.
\end{worked}

%---------------------------------------------------------------------
\begin{fourlenses}{Security Analytics}
\lensSEC{This chapter is the security lens in full. The distilled version: work
backwards from analyst capacity to a maximum false-positive rate; baseline every
entity against itself; run signatures and anomaly detection together; choose
features by manipulation cost; and treat the model as an asset that will itself be
attacked.}

\lensLA{Academic integrity analytics is structurally the same problem and inherits
the same arithmetic. Cheating is rare, so the base rate is low and precision will be
poor --- a plagiarism or exam-proctoring flag is \emph{evidence to investigate},
never a verdict. Students also adapt: publicise a detection rule and behaviour
changes to evade it, which is the adversarial dynamic in a setting where the
``adversary'' is someone you owe a duty of care. The ethical stakes here are higher
than in a SOC, and \chref{ethics} should be read alongside this.}

\lensPM{Security work has a distinctive project profile: success is invisible, and
the business case rests on losses that did not happen. Metrics that survive contact
with management are mean time to detect, mean time to respond, coverage of assets
monitored, and alerts per analyst --- not model accuracy. Budget for the fact that
detection engineering is continuous maintenance rather than a delivery with an end
date; a model shipped and left alone degrades faster here than anywhere else in this
book.}

\lensIOT{The two chapters meet directly. IoT fleets are a favoured target --- weak
credentials, unpatched firmware, plaintext protocols --- and a compromised device is
often first visible as \emph{anomalous telemetry}, meaning the maintenance detector
of \chref{iot} is also a security sensor. Conversely, telemetry can be forged, so
readings from a device are not trustworthy evidence about that device. Cross-device
and cross-physics consistency checks are the practical defence.}
\end{fourlenses}

%---------------------------------------------------------------------
\begin{lab}[Per-Entity Baselines and an Evasion Test]
\begin{lstlisting}[language=Python]
import pandas as pd, numpy as np

# --- 1. per-entity baselines: the core idea of this chapter ------------
def entity_features(logs, entity="src_ip", window="30D"):
    g = logs.set_index("ts").groupby(entity)
    f = pd.DataFrame(index=logs.index)

    for col in ["bytes_out", "distinct_ports", "failed_logins"]:
        cur  = g[col].transform(lambda s: s.rolling("1h").sum())
        base = g[col].transform(lambda s: s.rolling(window).mean())
        sd   = g[col].transform(lambda s: s.rolling(window).std())
        # z-score against the entity's OWN history, not the population's
        f[f"{col}_z"] = (cur - base) / sd.replace(0, np.nan)

    # "has this host ever done this before?" -- powerful and cheap
    f["novel_dest"] = g["dst_ip"].transform(
        lambda s: (~s.duplicated()).rolling(100).sum())
    return f

# --- 2. threshold from ANALYST CAPACITY, not from 0.5 ------------------
def threshold_for_capacity(scores, events_per_day, alerts_per_day=120):
    """Work backwards from what the SOC can actually investigate."""
    target_fpr = alerts_per_day / events_per_day
    return np.percentile(scores, 100 * (1 - target_fpr))

thr = threshold_for_capacity(scores, events_per_day=2_000_000, alerts_per_day=120)
print(f"threshold {thr:.4f} -> ~{(scores > thr).mean()*2e6:.0f} alerts/day")

# --- 3. TEST YOUR OWN MODEL FOR EVASION. This step is not optional. ----
def evasion_test(model, attack_rows, feature, factors=(0.9, 0.8, 0.6, 0.4)):
    """How much must an attacker reduce one feature to escape detection?"""
    base_rate = (model.predict_proba(attack_rows)[:, 1] > thr).mean()
    print(f"detection at full strength: {base_rate:.1%}")
    for k in factors:
        probe = attack_rows.copy()
        probe[feature] *= k                     # attacker throttles this feature
        rate = (model.predict_proba(probe)[:, 1] > thr).mean()
        print(f"  {feature} x{k}: detection {rate:.1%}")
        if rate < 0.2:
            print(f"  ^ EVADED by scaling one feature to {k}. Too fragile.")

# evasion_test(model, known_attacks, "distinct_ports_per_minute")

# --- 4. screen the training set before retraining (anti-poisoning) -----
from sklearn.ensemble import IsolationForest
screen = IsolationForest(contamination=0.01, random_state=0).fit(X_new)
clean = X_new[screen.predict(X_new) == 1]
print(f"withheld {len(X_new) - len(clean)} suspicious rows from retraining")
\end{lstlisting}

\textbf{Step 3 has no counterpart anywhere else in this book.} In every other
domain you test a model against held-out data; here you must also test it against
an opponent who is allowed to change the input.
\end{lab}

%---------------------------------------------------------------------
\begin{checkpoint}
\begin{tightnum}
  \item A detector with 99\% recall and a 1\% false-positive rate runs on 500{,}000
        daily events with 10 true attacks. How many alerts per day, and what
        precision? Is it deployable with four analysts?
  \item Why is \texttt{packets\_per\_second} a weaker feature than
        \texttt{fraction\_of\_connections\_to\_closed\_ports}?
  \item Automatic weekly retraining on production traffic is proposed. Name the
        attack this enables and one control that mitigates it.
\end{tightnum}
\end{checkpoint}

%---------------------------------------------------------------------
\begin{chaptersummary}
\begin{tight}
  \item Security violates the assumption underlying Part III: the data-generating
        process is an adversary who adapts to your model.
  \item Every source has a blind spot; encrypted traffic in particular leaves only
        metadata.
  \item Signature and anomaly detection are complements, not alternatives.
  \item Alert capacity, not accuracy, decides deployability. Work backwards from
        analysts per day to a maximum false-positive rate.
  \item Baseline each entity against its own history --- the single most productive
        feature idea in the domain.
  \item Models are attacked by evasion, poisoning, extraction and membership
        inference. The defences are mostly architectural: rate limits, data vetting,
        restricted outputs, human review.
  \item Choose features by manipulation cost as well as predictive power, and test
        your own model for evasion before an attacker does.
\end{tight}
\end{chaptersummary}

\begin{reviewq}
  \item \textbf{(MCQ)} An attack that corrupts the training data to weaken the next
        model is: (a)~evasion (b)~poisoning (c)~extraction (d)~membership inference.
  \item \textbf{(MCQ)} For a detector on 10 million daily events with a SOC capacity
        of 200 alerts/day, the maximum tolerable false-positive rate is about:
        (a)~$2\times10^{-5}$ (b)~$2\times10^{-3}$ (c)~0.02 (d)~0.2.
  \item \textbf{(Short)} Explain why ROC-AUC is a misleading headline metric for
        intrusion detection, and state what you would report instead.
  \item \textbf{(Short)} Define UEBA and explain why per-entity baselines outperform
        global thresholds.
  \item \textbf{(Short)} Describe membership inference and say why it is a privacy
        problem rather than a detection problem.
  \item \textbf{(Applied)} Design a detection system for compromised student
        accounts on a university VLE. Specify data sources, five features (at least
        three baselined per entity), the model, how you set the threshold, and two
        adversarial-ML risks with their mitigations.
  \item \textbf{(Applied)} Your organisation retrains its intrusion detector nightly
        on the previous day's traffic. Write the argument for and against this
        practice, and propose a concrete compromise procedure.
\end{reviewq}
